% ============ 第2章 §2.2 一般测度 (原书 pp.43-47) ============

\section{一般测度}\label{sec:2.2}
\index{0symbolsigma@$\sigma$-代数}

\index{H@荷重}

\index{C@测度}

\index{C@测度空间}










这里离开 Lebesgue 测度空间, 复习一般的测度空间. 本书中, 加权的情形与计数测度重要.

\subsection*{$\sigma$-代数}

调和分析中, 在配备 Lebesgue 测度的 Euclid 空间上分析是基本的, 但所考虑的集合不限于 Euclid 空间. 也有配备测度的距离空间的情形, 本书虽不涉及, 但也有 Riemann 流形的情形. 这里, 复习相当于为定义测度所需的定义域的 $\sigma$-代数, 以及广义化的测度的定义.

\begin{definition}[$\sigma$-代数]\label{def:2.10}
集合 $X$ 上的 $\sigma$-代数是指包含空集合、$X$ 本身, 并且对补集合、可数无限次的并及交的运算封闭的 $X$ 上的集合族.
\end{definition}

\begin{definition}[测度]\label{def:2.11}
设 $(X, \mathcal{M})$ 为可测空间. 即, $\mathcal{M}$ 是 $X$ 上的 $\sigma$-代数. $\mu$ 是测度是指

(1) $\mu: \mathcal{M} \to [0, \infty]$ 是集合函数.

(2) $\mu(\emptyset) = 0$.

(3) 当 $\{A_{j}\}_{j=1}^{\infty}$ 两两不交时,
\[
\mu\left(\sum_{j=1}^{\infty}A_{j}\right) = \sum_{j=1}^{\infty}\mu(A_{j})
\]
成立.
\end{definition}

\begin{definition}\label{def:2.12}
$(X, \mathcal{M}, \mu)$ 是测度空间是指

(1) $X$ 是集合.

(2) $\mathcal{M}$ 是 $X$ 上的 $\sigma$-代数.

(3) $\mu$ 是从 $\mathcal{M}$ 到 $[0, \infty]$ 的映射, 且是测度.
\end{definition}

作为 Lebesgue 测度以外的显著例子是下面这个. $2^{X}$ 是集合 $X$ 的全体子集构成的集合族.

\begin{definition}\label{def:2.13}
$\mu$ 是集合 $X$ 上的计数测度是指

(1) $\mu: 2^{X} \to [0, \infty]$ 是集合函数.

(2) $\mu(A) = \sharp A$.

成立.
\end{definition}

利用计数测度, 和作为积分被广义化. 作为从测度变换为别的测度的标准手法, 有加权的方法.

\begin{definition}\label{def:2.14}
设 $(X, \mathcal{M}, \mu)$ 为测度空间. $w: X \to (0, \infty)$ 为非负值可测函数时, 加权测度 $\nu = w\,\dd\mu$ 是指由
\[
\nu(E) = \int_{E}w(x)\dd\mu(x) \quad (E \in \mathcal{M})
\]
给出的测度.
\end{definition}

在这个定义中, $w$ 称为权.

\subsection*{几乎处处}

基于即使在测度 $0$ 的集合上函数不同也不影响积分值这一点, 给出下面的定义.

\begin{definition}\label{def:2.15}
某个性质 $P$ 在 $\mu$-a.e. (almost everywhere) 下成立, 是指性质 $P$ 不成立的点全体的 $\mu$-测度为 $0$. $\mu$ 为 Lebesgue 测度时简称为 a.e.. a.e. 称为``几乎处处''、``几乎所有'', 中文术语按副词、形容词分别使用.
\end{definition}

\subsection*{积分的定义等}

函数的可测性、单调收敛定理等也可以与 Lebesgue 测度的情形一样地定义. 这里给出其一例.

\begin{definition}\label{def:2.16}
设 $(X, \mathcal{M}, \mu)$ 为测度空间. $f: X \to [-\infty, \infty]$ 是 $\mathcal{M}$-可测 (或可测, $\mu$-可测) 是指
\[
\{f > \alpha\} = \{x \in X : f(x) > \alpha\} \in \mathcal{M} \quad (\alpha \in \mathbb{R})
\]
成立.
\end{definition}

\begin{theorem}[单调收敛定理]\label{thm:2.10}
设 $(X, \mathcal{M}, \mu)$ 为测度空间. 设 $E$ 为可测集合, $f_{1}, f_{2}, \ldots: E \to \overline{\mathbb{R}} = \mathbb{R} \cup \{\pm\infty\}$ 为可测函数. 若 $0 \leq f_{1} \leq f_{2} \leq \cdots$, 则
\[
\lim_{j \to \infty}\int_{E}f_{j}(x)\dd\mu(x) = \int_{E}\lim_{j \to \infty}f_{j}(x)\dd\mu(x)
\]
成立.
\end{theorem}

\begin{proof}
这里不作证明. 请参考适当的文献.
\end{proof}

\begin{theorem}[Fatou 引理]\label{thm:2.11}
设 $(X, \mathcal{M}, \mu)$ 为测度空间. 设 $E$ 为可测集合, $f_{1}, f_{2}, \ldots: E \to \overline{\mathbb{R}} = \mathbb{R} \cup \{\pm\infty\}$ 为可测函数. 若
\[
0 \leq f_{1}, f_{2}, \ldots,
\]
则
\[
\int_{E}\liminf_{j \to \infty}f_{j}(x)\dd\mu(x) \leq \liminf_{j \to \infty}\int_{E}f_{j}(x)\dd\mu(x)
\]
成立.
\end{theorem}

\begin{proof}
这里不作证明. 请参考适当的文献.
\end{proof}

\begin{theorem}[Lebesgue 收敛定理]\label{thm:2.12}
设 $(X, \mathcal{M}, \mu)$ 为测度空间. 设 $E$ 为可测集合, $f_{1}, f_{2}, \ldots: E \to \overline{\mathbb{R}} = \mathbb{R} \cup \{\pm\infty\}$ 为可测函数. 又设 (极限存在且) $\lim_{j \to \infty}f_{j}(x) = f(x)$ 在 $\mu$-a.e. 下成立. 又, 若可积函数 $g$ 对各自然数 $j$ 与 $x \in E$ 满足 $|f_{j}(x)| \leq g(x)$, 则
\[
\lim_{j \to \infty}\int_{E}f_{j}(x)\dd\mu(x) = \int_{E}\lim_{j \to \infty}f_{j}(x)\dd\mu(x)
\]
成立.
\end{theorem}

\begin{proof}
这里不作证明. 请参考适当的文献.
\end{proof}

\begin{theorem}\label{thm:2.13}
设 $(X, \mathcal{M}, \mu)$ 为测度空间. 若对 $a < t < b$, $f_{t}: E \to \overline{\mathbb{R}} = \mathbb{R} \cup \{\pm\infty\}$ 已被定义, 且满足以下条件, 则
\[
\frac{\dd}{\dd t}\int_{E}f_{t}(x)\dd\mu(x) = \int_{E}\partial_{t}f_{t}(x)\dd\mu(x)
\]
成立.

(1) 对 $t \in (a, b)$, $f_{t}$ 可积.

(2) $f_{t}(x)$ 关于 $\mu$-a.e. $x$ 可偏导.

(3) 存在可积函数 $g$, 当 $t \in (a, b)$ 时, 对 $\mu$-a.e. 有 $|\partial_{t}f_{t}(x)| \leq g(x)$ 成立.
\end{theorem}

\begin{proof}
这里不作证明. 请参考适当的文献.
\end{proof}

\subsection*{Fubini 定理}

作为特殊情形, 下面的概念重要. 特别是, $\sigma$-有限这一概念在考虑 Fubini 定理时重要. 在几乎所有情形下, 这个 $\sigma$-有限的条件都被满足.

\begin{definition}\label{def:2.17}
设 $(X, \mathcal{M}, \mu)$ 为测度空间.

(1) $\mu$ 是有限测度是指满足 $\mu(X) < \infty$.

(2) $\mu$ 是 $\sigma$-有限测度是指存在 $\mu$-测度有限的 $X_{1}, X_{2}, \ldots \in \mathcal{M}$ 使得
\[
X = \bigcup_{j=1}^{\infty}X_{j}.
\]
\end{definition}

单调收敛定理、Fatou 引理等特别地不需要对测度空间加限制, 但只有 Fubini 定理需要注意, 需要两个测度空间, 且这两个 $(X, \mathcal{M}, \mu)$ 与 $(Y, \mathcal{N}, \nu)$ 必须是 $\sigma$-有限的. 用 $\mathcal{M} \otimes \mathcal{N}$ 表示包含所有 $M \times N$ ($M \in \mathcal{M}$, $N \in \mathcal{N}$) 形状集合的 $X \times Y$ 上的最小 $\sigma$-代数. 又, 测度 $\mu \otimes \nu$ 满足 $\mu \otimes \nu(M \times N) = \mu(M)\nu(N)$, $M \in \mathcal{M}$, $N \in \mathcal{N}$.

\begin{theorem}[Fubini 定理]\label{thm:2.14}
设 $(X, \mathcal{M}, \mu)$ 与 $(Y, \mathcal{N}, \nu)$ 为 $\sigma$-有限的测度空间. $f$ 为非负值或复数值可积函数时,

(1) $f_{x} = f(x, \cdot)$ 关于 $\mu$-a.e. $x \in X$ 可积.

(2) 函数 $x \in X \mapsto \int_{Y}f(x, y)\dd\nu(y) \in [0, \infty]$ ($f$ 为非负值时), $x \in X \mapsto \int_{Y}f(x, y)\dd\nu(y) \in \mathbb{C}$ ($f$ 为复数值可积函数时) 是可积函数.

(3) 等式
\begin{align*}
\int_{X}\left(\int_{Y}f(x, y)\dd\nu(y)\right)\dd\mu(x) &= \int_{Y}\left(\int_{X}f(x, y)\dd\mu(x)\right)\dd\nu(y) \\
&= \int_{X \times Y}f(x, y)\dd\mu \otimes \nu(x, y)
\end{align*}
成立.
\end{theorem}

\begin{proof}
这里不作证明. 请参考适当的文献.
\end{proof}

\begin{example}\label{exa:2.2}
考虑配备计数测度的 $\mathbb{N}$.

(1) 函数 $f: \mathbb{N} \to \mathbb{C}$ 与复数列 $\{f(n)\}_{n=1}^{\infty}$ 具有完全相同的信息. 因此, $\mathbb{N}$ 上的函数与复数值数列被同一化.

(2) $L^{1}$-范数后述, 上述 $f$ 的 $L^{1}$-范数是 $\sum_{n=1}^{\infty}|f(n)| = \|\{f(n)\}_{n=1}^{\infty}\|_{\ell_{(\mathbb{N})}^{1}}$.

(3) 基于以上事实, 在这个测度空间上考虑 Fatou 引理时, 可如下表述. 对取值于 $[0, \infty]$ 的二重数列 $\{a_{m,n}\}_{m,n=1}^{\infty}$,
\[
\sum_{m=1}^{\infty}\liminf_{n \to \infty}a_{m,n} \leq \liminf_{n \to \infty}\sum_{m=1}^{\infty}a_{m,n}
\]
成立.
\end{example}

\begin{example}\label{exa:2.3}
对取值于 $[0, \infty]$ 的二重数列 $\{a_{m,n}\}_{m,n=1}^{\infty}$, 由 Fubini 定理, $\sum_{n=1}^{\infty}\sum_{m=1}^{\infty}a_{m,n} = \sum_{m=1}^{\infty}\sum_{n=1}^{\infty}a_{m,n}$ 成立.
\end{example}
